\documentclass[11pt,a4paper]{letter}
\usepackage[utf8]{inputenc}
\usepackage[ngerman]{babel}
\usepackage{amsmath, amssymb}
\usepackage{csquotes}
\usepackage{geometry}
\geometry{margin=2.5cm}

\begin{document}
\begin{letter}

Dipl.-Ing. (FH) Michael Czybor \\ Hauptstraße 78b \\ 4222 Langenstein, AT \\ deutscher Staatsbürger

\vspace{0.5cm}

An den Petitionsausschuss des Deutschen Bundestages \\ Platz der Republik 1 \\ 11011 Berlin

\begin{flushright}
	\today
\end{flushright}

\vspace{0.5cm}

\textbf{Betreff: Stellungnahme der Bundesregierung zu alternativen Kosmologien, Neubewertung der Forschungsmittelprioritäten und Prüfung der gesellschaftlichen Implikationen des Urknall-Paradigmas}

\vspace{0.5cm}

Sehr geehrte Damen und Herren,

hiermit reiche ich folgende Petition ein:

\vspace{0.3cm}

\noindent
\textbf{Petitionstext (Forderung)}

Der Deutsche Bundestag möge die Bundesregierung auffordern, dem \textbf{Ausschuss für Bildung, Forschung und Technikfolgenabschätzung} eine umfassende Stellungnahme zu folgenden Fragen vorzulegen:

\vspace{0.2cm}

\noindent
\textit{1. Zu alternativen Kosmologien:}

\begin{enumerate}
    \item Welche alternativen kosmologischen Modelle zur Urknall-Hypothese ($\Lambda$CDM) werden international diskutiert – insbesondere stationäre Kosmologien, Plasma-Kosmologie, fraktale Raumtheorien, Weber-Gravitation und die de Broglie-Bohm-Quantentheorie?
    
    \item Wie bewertet die Bundesregierung die Tatsache, dass nach Jahrzehnten intensiver Suche keine direkten Nachweise für dunkle Materie und dunkle Energie erbracht wurden – zwei zentrale Säulen des Urknall-Modells?
    
    \item Welche wissenschaftlichen und forschungspolitischen Konsequenzen zieht die Bundesregierung aus der Existenz konsistenter Alternativen, die ohne diese nicht nachweisbaren Entitäten auskommen?
\end{enumerate}

\noindent
\textit{2. Zur Neuausrichtung der Forschungsmittel:}

\begin{enumerate}
    \item Hält die Bundesregierung eine unabhängige, transparente Kosten-Nutzen-Analyse der deutschen Beteiligung an Urknall-orientierten Großforschungsprojekten (CERN, LISA, Euclid, Planck, LIGO/Virgo/KAGRA, Einstein-Teleskop, SKA, James-Webb-Teleskop etc.) für sinnvoll – unter Einbeziehung der Opportunitätskosten im Vergleich zu anderen Forschungsfeldern?
    
    \item Welche Prioritäten setzt die Bundesregierung bei der Verteilung von Forschungsmitteln zwischen spekulativer Grundlagenforschung (Urknall, dunkle Entitäten, Stringtheorie, Supersymmetrie) und anwendungsnaher Forschung mit gesellschaftlich unmittelbarem Nutzen (Fusionsenergie, Plasmaphysik, Kernfusion, erneuerbare Energien, Nachhaltigkeitstechnologien)?
    
    \item Wäre die Bundesregierung bereit, im Rahmen ihrer Gremienmitgliedschaften (CERN-Rat, ESA-Rat, ESO-Rat, SKA-Observatorium) eine Initiative zu starten, die eine vergleichende Bewertung alternativer kosmologischer Paradigmen auf internationaler Ebene anstößt?
\end{enumerate}

\noindent
\textit{3. Zu den gesellschaftlichen und weltanschaulichen Implikationen:}

\begin{enumerate}
    \item Ist der Bundesregierung die auffällige strukturelle Parallele zwischen dem expandierenden Universum der Urknall-Kosmologie und dem expansiven Wachstumszwang der kapitalistischen Wirtschaft bekannt – beide Systeme basieren auf unsichtbaren, nicht nachweisbaren Entitäten (dunkle Energie / dunkle Materie auf der einen Seite, \enquote{unsichtbare Hand} des Marktes / Externalisierung von Kosten auf der anderen Seite)?
    
    \item Welche Parallelen sieht die Bundesregierung zwischen dem postulierten Hitzetod des Universums (Big Rip, Big Crunch, Big Freeze) und den selbstzerstörerischen Tendenzen der gegenwärtigen Umwelt- und Klimapolitik – insbesondere der Förderung von Wachstum auf einem endlichen Planeten?
    
    \item Inwiefern ist die Urknall-Hypothese nach Ansicht der Bundesregierung nicht nur ein physikalisches Modell, sondern auch ein kulturelles Narrativ, das ein bestimmtes Wirtschafts- und Gesellschaftsmodell legitimiert – und welche Alternativen bieten ewige, stationäre Kosmologien (wie in der beigefügten Arbeit dargestellt) für ein Nachhaltigkeitsparadigma, eine Kreislaufwirtschaft und einen respektvollen Umgang mit der Natur?
    
    \item Hält die Bundesregierung eine öffentliche, interdisziplinäre Debatte im Deutschen Bundestag über diese Zusammenhänge (Physik, Kosmologie, Wirtschaftstheorie, Ökologie, Kulturwissenschaften) für erforderlich – und wenn ja, welche konkreten Schritte wird sie unternehmen, um eine solche Debatte anzustoßen (z. B. Anhörung von Sachverständigen, Einsetzung einer Enquete-Kommission, Bericht des Büros für Technikfolgen-Abschätzung)?
    
    \item Welche forschungspolitischen Konsequenzen zieht die Bundesregierung aus der grundsätzlichen Erkenntnis, dass die Entscheidung für oder gegen das Urknall-Paradigma keine rein physikalische Frage ist, sondern eine politische, kulturelle und weltanschauliche Grundsatzentscheidung mit erheblichen Auswirkungen auf das Selbstverständnis des Menschen und seinen Umgang mit der Natur?
\end{enumerate}

\vspace{0.3cm}

\noindent
\textbf{Begründung}

\vspace{0.2cm}

\noindent
\textit{1. Wissenschaftliche Grundlage}

Die Urknall-Hypothese basiert wesentlich auf den Postulaten der dunklen Materie und der dunklen Energie. Beide Entitäten wurden trotz massiver Forschungsanstrengungen (CERN, LUX-ZEPLIN, XENONnT, DarkSide, ADMX, Euclid, James-Webb-Teleskop, Planck-Satellit, LISA etc.) über Jahrzehnte hinweg nie direkt nachgewiesen. Die Suche nach WIMPs, Axionen, sterile Neutrinos, primordialen Schwarzen Löchern und anderen hypothetischen Teilchen blieb ergebnislos. Dies ist ein ernstzunehmendes wissenschaftliches Indiz dafür, dass das Urknall-Modell möglicherweise auf falschen Grundannahmen beruht.

Es existieren konsistente wissenschaftliche Alternativen, die ohne diese Postulate auskommen:

\begin{itemize}
    \item Die \textbf{de Broglie-Bohm-Theorie} ist eine anerkannte, deterministische Alternative zur Kopenhagener Deutung der Quantenmechanik. Sie vermeidet das Messproblem und die Nichtlokalität ist keine Paradoxie, sondern ein physikalisches Phänomen, das durch das Quantenpotential vermittelt wird.
    
    \item Die \textbf{Weber-Gravitation} (basierend auf den Arbeiten von Wilhelm Weber, wiederentdeckt und systematisch ausgearbeitet von André Koch Torres Assis) beschreibt Gravitation als geschwindigkeits- und beschleunigungsabhängige direkte Fernwirkung – ohne Raumzeitkrümmung, ohne dunkle Materie, ohne dunkle Energie.
    
    \item Die \textbf{$\Omega$-Theorie} (aus der Weber-Gravitation abgeleitet) formuliert Gravitation als Rotationsfeld, nicht als Raumzeitkrümmung. Sie erklärt Galaxienrotationskurven ohne dunkle Materie durch ein gravito-magnetisches Feld und anisotrope Trägheit.
    
    \item \textbf{Fraktale Raummodelle} (mit nicht-ganzzahliger Hausdorff-Dimension $D \approx 2,704$) ergeben zwangsläufig Skalengesetze, die viele kosmologische Beobachtungen ohne Expansion und ohne dunkle Komponenten erklären können – von Galaxienrotationskurven über die Hubble-Spannung bis hin zur großskaligen Struktur.
\end{itemize}

Die beigefügte Arbeit von Michael Czybor (\enquote{Die vereinheitlichte Theorie von Information, Raum und Gravitation} – IWT \& WDBT+) stellt ein geschlossenes alternatives Paradigma dar, das ohne Urknall, ohne dunkle Materie, ohne dunkle Energie, ohne Inflation und ohne Singularitäten auskommt. Sie ist axiomatisch geschlossen, parameterfrei, singularitätenfrei und macht testbare Vorhersagen (frequenzabhängige Lichtablenkung, Gravitationswellen-Dispersion, maximale Masse kompakter Objekte $M_{\text{BEC}} \approx 3 \times 10^6 M_\odot$, Neutrinomasse $m_\nu \approx 0,1\,\text{eV}/c^2$, fraktale Skalengesetze $\rho(r) \propto r^{-0,296}$). Diese Vorhersagen sind falsifizierbar und unterscheiden sich klar von den Vorhersagen des $\Lambda$CDM-Modells.

\vspace{0.2cm}

\noindent
\textit{2. Forschungspolitische Dimension}

Die Bundesregierung investiert massiv in die Fusionsforschung (2 Milliarden Euro bis 2029, Aufbau von drei Hubs: Magnetfusion in Greifswald, Laserfusion in Darmstadt, Brennstoffkreislauf in Jülich; \enquote{Fusionstalente}-Programm für Nachwuchsgruppen; überplanmäßige Mittel von 17,3 Millionen Euro zur Personalbindung). Dies ist ausdrücklich zu begrüßen, da die Fusionsenergie eine realistische Perspektive auf eine nachhaltige, saubere Energieversorgung bietet.

Zugleich fließen weiterhin jährlich Milliarden Euro in Urknall-orientierte Großforschung:
\begin{itemize}
    \item Deutsche Beiträge an CERN: ca. 250 Millionen Euro pro Jahr (etwa 20\% des LHC-Betriebsbudgets)
    \item ESA-Missionen mit deutscher Beteiligung: Euclid (dunkle Energie), LISA (Gravitationswellen), Planck (Hintergrundstrahlung), James-Webb-Teleskop
    \item Gravitationswellen-Observatorien: LIGO, Virgo, KAGRA, geplantes Einstein-Teleskop
    \item Radioteleskope: SKA (Square Kilometre Array)
    \item Dunkle-Materie-Detektoren: XENONnT, LUX-ZEPLIN, ADMX
\end{itemize}

Die wissenschaftliche Grundlage dieser Projekte – das Postulat nicht nachweisbarer dunkler Entitäten – ist, wie unter Punkt 1 dargelegt, zunehmend fragwürdig. Eine unabhängige, transparente Kosten-Nutzen-Analyse unter Einbeziehung alternativer kosmologischer Paradigmen ist daher nicht nur wünschenswert, sondern aus haushaltspolitischer Verantwortung heraus geboten.

\vspace{0.2cm}

\noindent
\textit{3. Die Parallele: Kosmologie als Spiegel der Gesellschaft}

Die Urknall-Kosmologie ist kein wertneutrales wissenschaftliches Modell. Sie transportiert ein Narrativ der ewigen Expansion aus einer Singularität, gestützt auf unsichtbare, nicht nachweisbare Entitäten – eine strukturell auffällige Parallele zum expansiven Wachstumszwang der kapitalistischen Wirtschaft, gestützt auf die \enquote{unsichtbare Hand} des Marktes und die Externalisierung von Kosten (Umwelt, soziale Ungleichheit, Ressourcenverbrauch).

Die postulierten Endzustände des Urknall-Universums (Big Rip, Big Crunch, Big Freeze, Hitzetod) finden ihre gesellschaftliche Entsprechung in den apokalyptischen Szenarien von Klimakollaps, Artensterben, Ressourcenerschöpfung und ökologischer Zerstörung. Die wissenschaftliche und politische Frage ist: Sind dies neutrale, wertfreie Vorhersagen – oder sich selbst erfüllende Prophezeiungen, die ein bestimmtes Wirtschafts- und Lebensmodell legitimieren, indem sie Alternativen als \enquote{unwissenschaftlich} diskreditieren?

Eine stationäre, ewige Kosmologie (wie in der beigefügten Arbeit dargestellt) wäre dagegen ein wissenschaftliches Fundament für Nachhaltigkeit, Kreislaufwirtschaft, Gleichgewicht und einen respektvollen Umgang mit der Natur. Die Entscheidung für oder gegen das Urknall-Paradigma ist daher keine rein physikalische Frage. Sie ist eine politische, kulturelle und weltanschauliche Grundsatzentscheidung mit erheblichen Auswirkungen auf das Selbstverständnis des Menschen und seinen Umgang mit der Natur.

Die beigefügte Arbeit von Michael Czybor zeigt einen konkreten, konsistenten, testbaren Weg für eine solche alternative Kosmologie auf – basierend auf der de Broglie-Bohm-Quantentheorie (anerkannte Alternative zur Kopenhagener Deutung), der Weber-Gravitation (direkte Fernwirkung, keine dunkle Materie) und der fraktalen Geometrie (die Natur ist fraktal, nicht glatt und dreidimensional).

\vspace{0.2cm}

\noindent
\textit{4. Zuständigkeit des Ausschusses}

Der Ausschuss für Bildung, Forschung und Technikfolgenabschätzung ist gemäß der Geschäftsordnung des Deutschen Bundestages zuständig für Grundsatzfragen der Forschungsförderung, für die Bewertung wissenschaftlicher Paradigmen sowie für die Abschätzung der gesellschaftlichen, technologischen und kulturellen Folgen von Forschungsentscheidungen. Die hier aufgeworfenen Fragen betreffen unmittelbar alle drei Bereiche.

\vspace{0.3cm}

\noindent
\textbf{Beigefügte Anlage}

\begin{itemize}
    \item Czybor, Michael: \enquote{Die vereinheitlichte Theorie von Information, Raum und Gravitation} (IWT \& WDBT+), Juni 2026 – vollständiges Dokument als Diskussionsgrundlage
\end{itemize}

\vspace{0.3cm}

\noindent
\textbf{Erklärung}

Ich versichere, dass ich die vorstehenden Angaben nach bestem Wissen und Gewissen gemacht habe. Mir ist bekannt, dass ich für die Richtigkeit der Angaben die Verantwortung trage.

\vspace{1.5cm}

\begin{flushleft}
    \rule{6cm}{0.4pt}
    
    \vspace{0.2cm}
    
    Langenstein/AT, \today\\Michael Czybor
\end{flushleft}

\end{letter}

\end{document}